\section{Проектиране на уеб-базираната платформа}

Настоящата глава представя проектирането на разработваната уеб-базирана платформа. След въвеждането на общия контекст и теоретичната рамка в предходната глава, тук фокусът се пренасочва към конкретните инженерни решения: изискванията към системата, нейната архитектурна организация и домейн моделът, който описва предметната област. Разгледани последователно, тези три аспекта формират проектантската основа на приложението и подготвят прехода към програмната реализация.

\subsection{Анализ на изискванията и потребителските сценарии}

След въвеждането на общия контекст и теоретичната рамка е необходимо системата да бъде описана от гледна точка на конкретните изисквания, които тя трябва да удовлетворява. Тази стъпка е важна, защото осигурява връзката между предметната област и реалната имплементация. Изискванията не са само списък от желани функции; те са критерий, спрямо който могат да бъдат оценени архитектурата, структурата на кода и степента, в която приложението изпълнява предназначението си.

В случая на онлайн магазин за компютърна и офис техника изискванията се формират както от типичните процеси в електронната търговия, така и от конкретния начин, по който проектът е структуриран в сорс кода. Анализът показва, че системата е ориентирана към сравнително малка, но функционално завършена среда, в която регистрираният потребител може да управлява своята информация, да въвежда фирми и продукти и да преминава през основния цикъл на покупка.

\subsubsection{Основни участници в системата}

В най-общ вид могат да бъдат разграничени три типа участници, които влияят върху поведението на системата.

\paragraph{Посетител без активна сесия}

Това е потребител, който достъпва приложението без да е логнат. В този режим системата следва да предоставя начална страница, страница за регистрация и страница за вход. Посетителят няма право да разглежда профилни данни, да добавя компании, да въвежда продукти или да използва количката. От гледна точка на реализирания код тази логика е постигната чрез проверки за наличие на атрибут \texttt{email} в сесията.

\paragraph{Потвърден регистриран потребител}

Това е основният активен участник в системата. След потвърждение на регистрацията чрез електронна поща и успешен вход той получава достъп до каталога, профила, списъка с потребители, фирмите, формите за добавяне на фирма и продукт, количката и историята на поръчките. Важно наблюдение е, че в текущия проект липсва отделна роля от тип администратор; много от действията, които в производствена среда биха били ограничени, тук са достъпни за всеки удостоверен потребител. Това е приемливо за учебен сценарий, но трябва да бъде отчетено като особеност на изискванията към настоящата имплементация.

\paragraph{Външни инфраструктурни компоненти}

Освен човешките потребители системата взаимодейства и с външни компоненти. На първо място това е MySQL базата данни, която съхранява постоянната информация. На второ място това е SMTP сървърът, използван за изпращане на електронно писмо при регистрация. На трето място може да се разглежда и браузърният JavaScript код, който използва REST крайна точка за извличане на максимална цена на продукт. Макар тези компоненти да не са „актьори“ в класически смисъл, те формират съществена част от оперативния контекст на приложението.

\subsubsection{Основни потребителски цели}

Анализът на предметната област позволява да бъдат формулирани няколко основни потребителски цели:

\begin{itemize}[leftmargin=*]
    \item създаване на акаунт и потвърждение на електронна поща;
    \item успешно удостоверяване и работа в персонална сесия;
    \item въвеждане и поддържане на фирмена информация, асоциирана с потребителя;
    \item въвеждане на продукти с категория, цена, количество и изображение;
    \item разглеждане на каталог и избор на подходящ артикул;
    \item добавяне на избран продукт в количка с желано количество;
    \item финализиране на поръчка с адрес за доставка;
    \item преглед на минали поръчки и история на потребителската активност.
\end{itemize}

Тези цели са практическият еквивалент на функционалните изисквания. Те показват как системата се възприема от потребителя не като набор от контролери и таблици, а като последователност от очаквани резултати.

\subsubsection{Функционални изисквания}

Функционалните изисквания определят какви операции системата трябва да поддържа. В настоящата разработка те могат да бъдат формулирани по следния начин.

\paragraph{Регистрация и потвърждение на акаунт}

Системата трябва да позволява създаване на нов потребителски профил чрез форма за регистрация. При успешно подадени данни следва да бъде създаден запис в базата, паролата да бъде съхранена в хеширан вид, а на посочения имейл да бъде изпратена връзка за потвърждение. До потвърждаване на акаунта достъпът до системата не трябва да бъде пълен. В текущата реализация това е постигнато чрез булево поле за потвърден статус и логика в потребителската услуга.

\paragraph{Вход и изход от системата}

След потвърждаване на акаунта потребителят трябва да може да влезе в системата с имейл и парола. При успешен вход в сесията се съхранява идентификаторът на потребителя под формата на електронна поща, което позволява контролерите да проверяват дали сесията е активна. Системата трябва да поддържа и коректен изход, при който сесията се инвалидира.

\paragraph{Управление на потребителски профил}

Удостовереният потребител трябва да може да преглежда и редактира основни лични данни. Това включва промяна на име, фамилия, имейл и парола. При смяна на паролата системата следва да изисква въвеждане на старата парола и да проверява коректността ѝ, преди да запази новата стойност в криптиран вид.

\paragraph{Управление на фирми}

Потребителят трябва да може да въвежда фирми, които впоследствие да се използват при добавяне на продукти. За всяка фирма се поддържат име, ЕИК/идентификатор, ДДС номер, адрес и признак за регистрация. В контекста на текущия проект фирмата е логически свързана с конкретен потребител и се използва като източник на собственост за продуктите.

\paragraph{Управление на продукти}

Системата трябва да позволява добавяне на нов продукт с наименование, цена, количество, категория, фирма и URL към изображение. Категорията следва да бъде избирана от предварително дефиниран набор от стойности, което улеснява контрола на данните и последващото филтриране в каталога. След добавяне продуктът трябва да бъде достъпен за визуализация в каталожната част на приложението.

\paragraph{Каталог и преглед на детайли}

Всеки удостоверен потребител трябва да може да разглежда каталога на наличните продукти. Каталогът трябва да поддържа филтриране по категории, а за всеки продукт да се визуализират име, категория, цена и изображение. При избор на конкретен артикул системата трябва да предоставя детайлен изглед, от който да бъде възможно добавяне към количката.

\paragraph{Количка и финализиране на поръчка}

След избор на продукт потребителят трябва да може да добави желано количество в количката, ако то е налично. При този процес системата трябва да проверява дали заявеното количество не надвишава текущата наличност. Количката трябва да визуализира избраните артикули, единичната и сумарната цена, както и да поддържа оформяне на поръчка чрез адрес за доставка и калкулация на крайна стойност.

\paragraph{История на поръчките}

След приключване на покупката системата трябва да съхрани исторически запис за направената поръчка, включително дата, адрес и обща стойност. Потребителят трябва да може да разглежда тези записи в отделен екран и при необходимост да изчиства натрупаната история.

\paragraph{Локализация и допълнителен REST интерфейс}

Системата трябва да поддържа базова смяна на езика между английски и български за част от интерфейса чрез локализирани съобщения. Освен това трябва да предлага поне една примерна REST крайна точка, която връща агрегирана информация в машинно четим формат. В текущия проект това е информацията за максималната цена на продукт.

\begin{table}[H]
    \centering
    \caption{Обобщение на основните функционални изисквания}
    \label{tab:functional-requirements}
    \begin{tabularx}{\textwidth}{L{3.0cm}YY}
        \toprule
        Идентификатор & Изискване & Реализация в проекта \\
        \midrule
        FR-1 & Регистрация на потребител & Форма за регистрация, запис в базата, хеширане на парола \\
        FR-2 & Потвърждение на акаунт & Имейл с линк и поле \texttt{isConfirmed} \\
        FR-3 & Вход и изход & Проверка на парола и управление на \texttt{HttpSession} \\
        FR-4 & Редакция на профил & Форма за профил и проверка на стара парола \\
        FR-5 & Добавяне на фирма & Отделен контролер и binding модел за фирмени данни \\
        FR-6 & Добавяне на продукт & Категории, фирма, количество, цена и URL изображение \\
        FR-7 & Разглеждане на каталог & Категорийни маршрути и Thymeleaf визуализация \\
        FR-8 & Добавяне в количка & Детайлен изглед, количество и проверка за наличност \\
        FR-9 & Финализиране на поръчка & Адрес, обща цена и създаване на история \\
        FR-10 & История на поръчките & Екран с исторически записи и операция за изчистване \\
        FR-11 & Локализация & \texttt{messages.properties}, \texttt{messages\_bg.properties}, interceptor \\
        FR-12 & REST услуга & Крайна точка за максимална цена на продукт \\
        \bottomrule
    \end{tabularx}
\end{table}

Таблица \ref{tab:functional-requirements} показва, че функционалностите на системата следват логична последователност от регистрация и вход към управление на данни, покупка и проследяване на резултата от тази покупка. Това е именно цикълът, който определя проекта като завършено макар и компактно електронно търговско приложение.

\subsubsection{Нефункционални изисквания}

Наред с конкретните функции, системата следва да отговаря и на редица нефункционални изисквания. Те не описват какво прави приложението, а как следва да го прави.

\paragraph{Поддържаемост}

Кодът трябва да бъде организиран така, че логиката на системата да бъде лесна за проследяване. Това предполага разделяне по пакети, отделяне на контролери, услуги и repository компоненти, както и използване на отделни модели за вход, вътрешен трансфер и визуализация. Поддържаемостта е особено важна в дипломна разработка, тъй като документът трябва да може ясно да обясни структурата на проекта.

\paragraph{Четим и последователен потребителски интерфейс}

Системата трябва да предлага интерфейс, който е достатъчно интуитивен за базови потребителски операции. Макар дизайнът да не е водещ изследователски фокус, от приложението се очаква да поддържа последователна навигация, разбираеми форми и ясни маршрути към основните действия.

\paragraph{Коректност на входните данни}

Всички основни форми следва да бъдат защитени чрез валидация, така че очевидно невалидни стойности да не достигат до persistence слоя. В разглеждания проект това се реализира чрез анотации и обработка на резултата от bind операцията.

\paragraph{Устойчиво съхранение на информацията}

Приложението трябва да използва релационна база данни, за да гарантира, че потребители, фирми, продукти, количка и история се съхраняват устойчиво и са достъпни между отделни изпълнения на системата. Използването на MySQL и Hibernate осигурява именно този тип устойчивост.

\paragraph{Елементарна сигурност}

Необходимо е паролите да бъдат съхранявани хеширано, а достъпът до определени страници да бъде ограничен до удостоверени потребители. Макар това да не реализира пълноценна корпоративна схема за сигурност, то е минимално необходимо условие, за да бъде системата функционално и методологично коректна.

\paragraph{Разширяемост}

Архитектурата следва да позволява добавяне на нови категории, нови екрани и нови услуги без фундаментално пренаписване на вече изградената основа. Именно поради това изборът на многослоен модел и Spring компоненти има не само настоящ, но и стратегически смисъл.

\subsubsection{Ключови потребителски сценарии}

За по-голяма яснота функционалните изисквания могат да бъдат представени и като сценарии. Това е полезно, защото показва не просто кои функции съществуват, а как те се комбинират в завършен работен поток.

\paragraph{Сценарий 1: Регистрация и първи вход}

Посетителят отваря началната страница и избира регистрация. Той попълва име, фамилия, имейл и парола. След успешно изпращане системата създава нов потребител, хешира паролата и изпраща линк за потвърждение. След активиране на акаунта потребителят се връща към формата за вход, въвежда данните си и системата създава активна сесия.

\paragraph{Сценарий 2: Добавяне на фирма и продукт}

След успешен вход потребителят отваря екрана за добавяне на фирма и въвежда необходимите атрибути. След това преминава към формата за нов продукт, избира една от наличните категории, посочва фирмата, задава цена, количество и URL на изображение. След запис продуктът става част от каталога.

\paragraph{Сценарий 3: Разглеждане на каталог и покупка}

Потребителят разглежда каталога, филтрира по категория и избира конкретен продукт. Отваря детайлния екран, въвежда желано количество и го добавя в количката. Ако количеството е налично, системата актуализира наличността и създава или обновява запис в количката. След това потребителят отваря количката, въвежда адрес за доставка и потвърждава поръчката. Системата създава исторически запис и изчиства активната количка.

\paragraph{Сценарий 4: Преглед на история и профил}

Във всеки момент удостовереният потребител може да отвори профила си, да редактира лични данни и да провери асоциираните компании. Той може да отвори и историята на поръчките, където да види натрупаните записи и при необходимост да ги изчисти.

\subsubsection{Ограничения, произтичащи от текущата постановка}

Анализът на изискванията разкрива и някои характерни ограничения. Липсата на отделни роли означава, че системата не разграничава административни от потребителски действия. Локализацията е частична и не обхваща изцяло всички текстове и категории. Някои операции, които в production среда биха били защитени чрез по-строга политика, тук са реализирани с по-прост модел. Същевременно тези ограничения не обезсмислят проекта; напротив, те очертават ясна граница между учебна система и пълноценна комерсиална платформа.

\subsubsection{Изводи}

Изискванията към системата показват, че проектът е замислен като компактно, но завършено уеб приложение, покриващо най-важните процеси на един онлайн магазин. Формулираните функционални и нефункционални критерии създават здрава основа за анализа на архитектурата и кода в следващите раздели. Особено важно е, че изискванията са достатъчно конкретни, за да позволят последваща верификация, и достатъчно реалистични, за да направят системата представителна за практическо приложение в областта на електронната търговия.

\subsection{Архитектура на системата}

След формулирането на изискванията естествената следваща стъпка е анализ на архитектурната организация на приложението. Архитектурата е гръбнакът на всяка нетривиална система, защото определя как отделните отговорности са разпределени, как се движат данните между модулите и до каква степен проектът може да се поддържа и развива във времето. В настоящата разработка архитектурният избор е насочен към ясно разделение между уеб слой, бизнес логика, достъп до данни, конфигурация и визуализация. Този раздел описва именно това разпределение и го свързва с реалната структура на сорс кода.

\subsubsection{Обща архитектурна постановка}

Системата е реализирана като монолитно уеб приложение върху Spring Boot. Терминът „монолитно“ тук не следва да се възприема като негативен. В контекста на настоящия проект той означава, че всички основни функционални компоненти -- уеб маршрутизация, бизнес логика, persistence слой, HTML шаблони и конфигурация -- се намират в един проект и се стартират в един процес. Това решение е подходящо поради няколко причини.

Първо, домейнът на системата е относително компактен. Няма независими бизнес домейни, които да изискват отделни услуги или самостоятелни deployment единици. Второ, учебната цел на проекта е да демонстрира завършена, но прозрачна архитектура, а не да усложнява излишно инфраструктурния модел. Трето, Spring Boot предоставя достатъчно добра организационна рамка, така че монолитното приложение да остане вътрешно добре структурирано.

Архитектурната схема може да бъде представена като съвкупност от следните слоеве:

\begin{itemize}[leftmargin=*]
    \item слой за HTTP комуникация и управление на заявки;
    \item слой за бизнес операции и домейн логика;
    \item слой за достъп до данните и persistence;
    \item слой за представяне чрез HTML шаблони;
    \item конфигурационен слой за bean компоненти, локализация и инфраструктурни настройки.
\end{itemize}

\begin{figure}[H]
    \centering
    \begin{tikzpicture}
        \node[appblock] (client) {Клиентски браузър};
        \node[appblock, below=1.3cm of client] (web) {Web слой\\контролери и маршрути};
        \node[appblock, below=1.3cm of web] (service) {Service слой\\бизнес операции};
        \node[appblock, below=1.3cm of service] (repo) {Repository слой\\JPA интерфейси};
        \node[appblock, below=1.3cm of repo] (db) {MySQL};
        \node[appblock, right=4.2cm of web] (view) {Thymeleaf шаблони};
        \node[appblock, left=4.2cm of web] (rest) {REST представяне};
        \node[appaccent, right=4.2cm of service] (config) {Конфигурация\\Bean-и и i18n};

        \draw[appline] (client) -- (web);
        \draw[appline] (web) -- (service);
        \draw[appline] (service) -- (repo);
        \draw[appline] (repo) -- (db);
        \draw[appline] (web) -- (view);
        \draw[appline] (view) |- (client);
        \draw[appline] (web) -- (rest);
        \draw[appline] (rest) |- (client);
        \draw[appline] (config) -- (service);
        \draw[appline] (config) -- (web);
    \end{tikzpicture}
    \caption{Високо ниво на архитектурата на приложението}
    \label{fig:high-level-architecture}
\end{figure}

Фигура \ref{fig:high-level-architecture} показва, че приложението следва централен уеб поток, но може да предоставя резултати както чрез HTML визуализация, така и чрез JSON представяне. Това е важен белег за архитектурна гъвкавост, въпреки че системата е доминиращо сървърно рендирана.

\subsubsection{Пакетна организация на сорс кода}

Основният Java код е разположен в директорията \texttt{src/main/java/com/example/diplomenproekt}. Вътре в нея проектът е разделен на няколко логически под-пакета. Тази организация не е просто въпрос на стил; тя изразява пряко архитектурното разделение на отговорностите.

\begin{table}[H]
    \centering
    \caption{Пакети в Java слоя и тяхното предназначение}
    \label{tab:package-organization}
    \begin{tabularx}{\textwidth}{L{3.5cm}Y}
        \toprule
        Пакет & Основно предназначение \\
        \midrule
        \texttt{config} & Съдържа конфигурационни класове за bean дефиниции и локализация \\
        \texttt{models.entities} & Описва JPA entity обектите на домейна \\
        \texttt{models.bindingModel} & Съдържа модели за приемане и валидация на входни форми \\
        \texttt{models.serviceModel} & Описва междинни модели за пренос на данни към service слоя \\
        \texttt{models.viewModel} & Съдържа модели, предназначени за визуализация в шаблоните или API \\
        \texttt{repositories} & Дефинира интерфейсите за достъп до базата данни чрез Spring Data JPA \\
        \texttt{services} & Съдържа service интерфейси, описващи бизнес възможностите \\
        \texttt{services.impl} & Реализира бизнес логиката и координацията между компонентите \\
        \texttt{web} & Съдържа MVC и REST контролери за входящите HTTP заявки \\
        \bottomrule
    \end{tabularx}
\end{table}

От таблица \ref{tab:package-organization} се вижда, че дори при сравнително малък проект е постигната смислена пакетна сегрегация. Това е едно от силните качества на разработката, защото улеснява навигацията, поддръжката и аналитичното описание на системата.

\subsubsection{Входна точка и стартиране на приложението}

Входната точка на системата е класът \texttt{DiplomenProektApplication}, маркиран с анотация \texttt{@SpringBootApplication}. Тази анотация обединява механизми за component scanning, auto-configuration и bootstrapping на приложението \cite{springboot}. На практика това означава, че при стартиране Spring открива контролерите, услугите, repository интерфейсите и конфигурационните класове, създава необходимите bean инстанции и подготвя средата за обслужване на заявки.

От архитектурна гледна точка е важно, че входният клас не съдържа бизнес логика. Той изпълнява единствено ролята на стартова точка. Това е добра практика, защото предотвратява смесването на инфраструктурна и домейн логика още на най-високо ниво на приложението.

\subsubsection{Конфигурационен слой}

В пакета \texttt{config} са разположени два съществени класа: \texttt{ApplicationBeanConfiguration} и \texttt{I18nConfig}. Първият дефинира bean компоненти за \texttt{ModelMapper} и \texttt{PasswordEncoder}. Тук архитектурната идея е ясна: компоненти с инфраструктурна роля, използвани на множество места в системата, трябва да се управляват централизирано от Spring контейнера, а не да се създават ръчно във всеки клас.

\texttt{I18nConfig} от своя страна настройва \texttt{LocaleResolver} и \texttt{LocaleChangeInterceptor}, с което създава базов механизъм за локализация на интерфейса. Това показва, че архитектурата не е ограничена единствено до CRUD логика, а обхваща и допълнителни крос-срезови характеристики на приложението.

Централизираната конфигурация има и друго предимство: когато даден инфраструктурен избор трябва да бъде променен, това става на едно определено място. Така например подмяна на password encoder или промяна на локализационния механизъм не изисква търсене и редакция в множество несвързани класове.

\subsubsection{Уеб слой: контролери и маршрути}

Уеб слоят е реализиран в пакета \texttt{web}. В него се намират контролерите \texttt{HomeController}, \texttt{UserController}, \texttt{CompanyController}, \texttt{ProductController} и \texttt{MaxPriceRestController}. Разделянето им по функционални области е логично и съответства на основните домейн сценарии на системата.

\texttt{HomeController} обслужва началната точка на приложението и реализира проста логика за пренасочване към каталога при активна сесия. \texttt{UserController} е най-богатият по отношение на потоци компонент, защото управлява регистрацията, входа, изхода, профила, списъка с потребители и потвърждението на електронна поща. \texttt{CompanyController} обслужва формите за въвеждане на фирма. \texttt{ProductController} управлява продуктовия каталог, добавянето на продукт, детайлния изглед, количката, checkout процеса и историята на поръчките. \texttt{MaxPriceRestController} предоставя отделен JSON интерфейс.

От архитектурна гледна точка силна страна на проекта е, че контролерите до голяма степен запазват ролята си на координатори, а не на място за тежка бизнес логика. Те валидират достъпа до маршрута, проверяват резултата от binding операцията, извикват съответна услуга и определят към коя view да премине изпълнението. Това следва добрата практика контролерът да бъде „тънък“ и ориентиран към уеб потока, а не към самото бизнес правило.

\subsubsection{Service слой: сърце на бизнес логиката}

Service слоят е организиран чрез интерфейси и имплементации. Това дава две важни предимства. Първо, архитектурно се разграничават договорът и конкретната реализация. Второ, системата става по-подходяща за бъдещо тестване и подмяна на конкретни имплементации.

В настоящия проект могат да бъдат разграничени следните основни service области:

\begin{itemize}[leftmargin=*]
    \item потребителски операции -- регистрация, вход, потвърждение, редакция на профил, изчисляване на обща цена в количката;
    \item управление на фирми -- добавяне и извличане на фирми за конкретен потребител;
    \item управление на продукти -- добавяне, категоризация, търсене по идентификатор, проверка за наличност, добавяне в количка;
    \item количка -- съхранение и извличане на записи за избрани продукти;
    \item история -- финализиране на поръчка, изчистване и извличане на минали покупки.
\end{itemize}

Service слоят е ключов за архитектурата, защото именно там се вземат решения за последователността на операциите. Например при checkout процеса е необходимо не просто да се създаде исторически запис, а да се координира обновяването на количката, историята и продуктите. Този тип действия са типичен пример за логика, която не принадлежи нито на контролера, нито на repository интерфейса, а на самия бизнес слой.

\subsubsection{Persistence слой и repository компоненти}

Persistence слоят е реализиран чрез Spring Data JPA repository интерфейси: \texttt{UserRepository}, \texttt{CompanyRepository}, \texttt{ProductRepository}, \texttt{ShoppingCartRepository} и \texttt{HistoryRepository}. Тяхната основна функция е да предоставят стандартизиран и сравнително декларативен достъп до данните.

Наличието на методи като \texttt{findByEmail}, \texttt{getAllByUser\_Id}, \texttt{findAllByUser\_Id} и \texttt{deleteAllByUser\_Id} показва типичния за Spring Data подход: част от необходимата query логика се изразява директно чрез имената на методите, без ръчно писане на SQL. Това спестява boilerplate код и прави част от persistence слоя по-компактен.

От архитектурна перспектива repository слоят е правилно ограничен до операции по извличане, запис и изтриване. Той не съдържа самостоятелна бизнес логика, а служи като механизъм, чрез който service слоят реализира домейн правилата. Именно това разделение е един от важните признаци за архитектурна последователност.

\subsubsection{Модели за трансфер на данни между слоевете}

Интересна архитектурна особеност на проекта е използването на няколко типа модели:

\begin{itemize}[leftmargin=*]
    \item entity модели за persistence слоя;
    \item binding модели за входни форми;
    \item service модели за междинен трансфер;
    \item view модели за визуализация и API представяне.
\end{itemize}

Този подход съответства на добрите практики за ограничаване на директната зависимост между базата данни и потребителския интерфейс. Например не е необходимо HTML шаблоните да получават винаги самите entity обекти; вместо това те могат да работят с view модели, съдържащи само необходимите за визуализация полета. Аналогично, при приемане на данни от форма е по-удачно да се използва binding модел с валидационни анотации, отколкото директно persistence обектът.

Използването на \texttt{ModelMapper} подсказва стремеж към намаляване на ръчните преобразувания между тези представяния. Това е полезно от инженерна гледна точка, но не отменя необходимостта от внимателно моделиране на самите класове и полета.

\subsubsection{Представящ слой: шаблони и статични ресурси}

Визуалният слой на приложението е разположен в \texttt{src/main/resources/templates} и \texttt{src/main/resources/static}. Шаблоните отразяват основните потребителски екрани: начална страница, регистрация, вход, профил, добавяне на фирма, добавяне на продукт, каталог, детайли, количка, история и страници за грешки. Статичните ресурси включват CSS и JavaScript файлове, както и изображения.

Архитектурното значение на този слой се състои в това, че той остава относително изолиран от persistence логиката. Шаблоните работят с вече подготвени атрибути от модела и не съдържат сложни бизнес решения. Това е важен белег за дисциплинирано разделение на отговорностите.

\subsubsection{Архитектурен поток на ключови операции}

За да бъде по-ясно как слоевете си взаимодействат, е полезно да се разгледа общ модел на няколко основни операции.

\paragraph{Регистрация}

При регистрация браузърът изпраща POST заявка към потребителския контролер. Контролерът приема binding модел и го предава към потребителската услуга. Услугата проверява дали вече съществува потребител с този имейл, създава entity обект, хешира паролата, записва новия потребител чрез repository слоя и изпраща потвърждаващ имейл чрез инфраструктурен компонент. След това контролерът връща подходящо пренасочване.

\paragraph{Добавяне на продукт}

Контролерът за продукти валидира входната форма и предава данните към \texttt{ProductsService}. Услугата намира избраната фирма, създава нов продукт, преобразува категорията от низ към \texttt{CategoryEnum} и записва обекта в базата. По този начин контролерът остава концентриран върху маршрута и валидацията, а service слоят поема домейн решението.

\paragraph{Checkout}

Checkout сценарият е най-показателен за архитектурната координация. Контролерът приема адрес и обща цена, валидира входа и извиква \texttt{HistoryService}. Тази услуга създава запис в историята, асоциира го с потребителя, нулира или изчиства данните за текущата количка и премахва продукти с нулева наличност. Тук ясно се вижда защо service слоят е необходим: операцията изисква координация между няколко persistence компонента и не може да бъде сведена до една repository заявка.

\subsubsection{Силни страни на архитектурното решение}

От направения анализ могат да бъдат изведени няколко силни страни на архитектурата:

\begin{itemize}[leftmargin=*]
    \item ясна пакетна организация, която отразява логическите роли на класовете;
    \item разграничение между контролери, услуги и repository интерфейси;
    \item използване на отделни модели за различни етапи на обмен на данни;
    \item централизация на инфраструктурни компоненти чрез конфигурационни bean дефиниции;
    \item наличие както на HTML, така и на примерен REST интерфейс.
\end{itemize}

Тези характеристики правят системата подходяща не само за изпълнение на текущата функционалност, но и за последващ анализ и разширение.

\subsubsection{Архитектурни компромиси и ограничения}

Наред със силните страни е необходимо да бъдат отчетени и определени компромиси. Част от проверките за достъп са реализирани директно в контролерите чрез условие върху сесията, вместо чрез централизирана security архитектура. Това прави системата по-лесна за разбиране, но по-слаба откъм формална сигурност. Някои операции в service слоя показват силна зависимост от конкретната структура на данните, което е очаквано за малък проект, но при по-голямо приложение би следвало да бъде допълнително абстрахирано.

Допълнително, макар архитектурата да е многослойна, тя все още е тясно обвързана с конкретната реализация на сесиен достъп и локална конфигурация. Това е разумно за учебен сценарий, но следва да бъде адресирано при бъдещо развитие към по-зряла среда.

\subsubsection{Архитектурно обобщение}

Архитектурата на разглежданата система представлява добре подредена монолитна Spring Boot реализация с ясно разграничени слоеве и пакетна структура. Използваният модел е напълно адекватен за домейн като онлайн магазин с ограничен, но реалистичен обхват. Той позволява логическа проследимост на данните и операциите, създава добра основа за поддръжка и надграждане и подчертава учебно-инженерната стойност на разработката.

\subsection{Домейн модел и модел на данните}

Домейн моделът е централната абстракция, чрез която софтуерната система представя предметната област. В разглеждания проект той описва основните участници и артефакти в процеса на електронна търговия: потребител, фирма, продукт, количка и история на поръчки. Макар приложението да не реализира сложна бизнес семантика от enterprise клас, домейн моделът е достатъчно богат, за да наложи ясно картографиране към релационна база данни и разделяне на различни типове модели за обработка и визуализация на данни.

Настоящият раздел разглежда едновременно логическото съдържание на домейн обектите и начина, по който те са реализирани чрез JPA entity класове, binding модели, service модели и view модели. По този начин се проследява целият път на информацията: от въвеждането ѝ във форма, през вътрешното ѝ представяне в системата, до записа ѝ в базата и повторната ѝ визуализация.

\subsubsection{Основни домейн понятия}

В предметната област на електронния магазин могат да бъдат идентифицирани няколко базови понятия.

\begin{itemize}[leftmargin=*]
    \item \textbf{Потребител} -- физическо лице, което се регистрира, влиза в системата, поддържа профил, добавя фирми, въвежда продукти, използва количка и има история на поръчки.
    \item \textbf{Фирма} -- обект, който представя търговска или доставна организация, асоциирана с конкретен потребител и използвана при добавяне на продукти.
    \item \textbf{Продукт} -- артикул от каталога, описан чрез име, цена, наличност, изображение и категория.
    \item \textbf{Запис в количка} -- конкретно желание за покупка на определен продукт в дадено количество.
    \item \textbf{Исторически запис} -- обобщена информация за приключена поръчка, включваща адрес, дата и цена.
\end{itemize}

Тези понятия образуват ядрото на домейна и определят по-голямата част от класовете в \texttt{models.entities}. От гледна точка на \textit{domain-driven} мисленето \cite{evans}, те могат да бъдат разглеждани като основните бизнес обекти, около които се организират и останалите механизми в приложението.

\subsubsection{Entity моделът в JPA слоя}

Persistence слоят използва JPA entity класове, които наследяват \texttt{BaseEntity}. Макар в текущата работа да не е необходимо да се навлиза в детайли на базовия клас, неговото наличие подсказва стандартизирано централизиране на идентификатора за всички трайни обекти. Това е типична практика и намалява дублирането на общи persistence характеристики.

\paragraph{Потребител}

Класът \texttt{User} е основният aggregate обект в системата. Той съдържа име, фамилия, парола, признак дали акаунтът е потвърден и електронна поща. Освен това поддържа връзки към фирмите на потребителя, към записите в количката и към историята на поръчките. Така именно потребителят е логическият център на множество операции в приложението.

Особен интерес представлява полето за електронна поща. То е маркирано като уникално и допълнително има регулярен израз на ниво entity, което подсказва стремеж към защита на едно от най-важните идентифициращи полета. Паролата не се валидира на това ниво, а се обработва и хешира чрез service логика. Булевото поле за потвърждение е част от модела на регистрация и контролира допустимостта на входа.

\paragraph{Фирма}

Класът \texttt{Company} описва фирмена информация и включва атрибути като име, \texttt{mol}, \texttt{vat}, адрес и признак за регистрация. В системата фирмата принадлежи на конкретен потребител чрез връзка \texttt{@ManyToOne}. Освен това една фирма има множество продукти чрез \texttt{@OneToMany(mappedBy = "company")}. От логическа гледна точка това е разумен модел: фирмата изпълнява ролята на притежател или източник на предлаганите продукти.

Следва да се отбележи, че името на колекцията в класа \texttt{Company} е \texttt{companies}, въпреки че тя реално съдържа продукти. Това е незначително от функционална гледна точка, но е показателно за важността на консистентното именуване в домейн модела. Подобни дребни несъответствия не чупят системата, но намаляват концептуалната яснота и следва да бъдат разглеждани като потенциален обект на бъдещо рефакториране.

\paragraph{Продукт}

Класът \texttt{Product} е основният каталожен обект в приложението. Той съдържа име, цена, количество, URL към изображение, категория и връзка към фирма. Категорията е реализирана като изброим тип \texttt{CategoryEnum}, което е добър архитектурен избор. По този начин наборът от допустими категории е изрично описан в кода и не зависи от свободен текст или произволни стойности, въведени от потребителя.

Валидирането на цена и количество е реализирано чрез \texttt{@Min(0)}, което формализира базовото бизнес правило, че продуктите не могат да имат отрицателни стойности за тези полета. От бизнес гледна точка продуктът е едновременно обект на визуализация и на транзакционно взаимодействие, защото участва в количката и се променя при checkout процеса.

\paragraph{Запис в количка}

Класът \texttt{ShoppingCartEntity} моделира преходно търговско състояние: конкретен потребител е избрал конкретен продукт в конкретно количество. Той съдържа само три основни атрибута -- желано количество, продукт и потребител. Макар моделът да е опростен, той е достатъчен за целите на проекта и илюстрира идеята, че количката не е просто списък от продукти, а списък от връзки „потребител--продукт--брой“.

Интересен детайл е, че продуктът в количката е реализиран с \texttt{@OneToOne}, а потребителят с \texttt{@ManyToOne}. Това позволява множество записи в количката да принадлежат на един потребител, но ограничава логическата връзка към конкретен продукт. В production сценарий често би се използвала и различна схема за моделиране на поръчковите позиции, но за настоящия проект този модел е достатъчен, за да представи основната идея.

\paragraph{История на поръчките}

Класът \texttt{HistoryEntity} представя обобщен запис за реализирана поръчка. В него се съхраняват адрес, дата, цена и връзка към потребител. Това е умишлено опростен модел. Вместо поръчката да се разглежда като самостоятелен aggregate с отделни позиции, статуси и платежни детайли, тук тя е сведена до исторически запис за случило се checkout събитие. Подобен подход е напълно оправдан в учебно-приложен контекст, тъй като позволява системата да покаже завършен цикъл на покупка без да навлиза в сложността на пълномащабен order management модел.

\begin{figure}[H]
    \centering
    \begin{tikzpicture}
        \node[appblock] (user) {User};
        \node[appblock, below left=2.2cm and 2.3cm of user] (company) {Company};
        \node[appblock, below right=2.2cm and 2.3cm of user] (history) {HistoryEntity};
        \node[appblock, below=2.2cm of user] (cart) {ShoppingCartEntity};
        \node[appblock, below=2.2cm of company] (product) {Product};

        \draw[appline] (company) -- node[left]{Many-to-One} (user);
        \draw[appline] (history) -- node[right]{Many-to-One} (user);
        \draw[appline] (cart) -- node[right]{Many-to-One} (user);
        \draw[appline] (product) -- node[left]{Many-to-One} (company);
        \draw[appline] (cart) -- node[below]{One-to-One} (product);
    \end{tikzpicture}
    \caption{Обобщена схема на връзките между основните entity обекти}
    \label{fig:entity-relationships}
\end{figure}

Фигура \ref{fig:entity-relationships} представя опростено отношенията между основните трайни обекти. От нея ясно се вижда, че потребителят е централният домейн елемент, към който се асоциират фирмите, количката и историята, докато продуктите са свързани с фирмите и се използват в количката.

\subsubsection{Атрибутен анализ на основните entity класове}

За по-систематичен преглед е полезно ключовите класове да бъдат сравнени по техните атрибути и роли в системата.

\begin{longtable}{L{2.7cm}L{3.8cm}L{7.3cm}}
\caption{Атрибути и роля на основните entity класове}\label{tab:entity-attributes}\\
\toprule
Клас & Основни атрибути & Роля в системата \\
\midrule
\endfirsthead
\toprule
Клас & Основни атрибути & Роля в системата \\
\midrule
\endhead
\texttt{User} & \texttt{firstName}, \texttt{lastName}, \texttt{password}, \texttt{email}, \texttt{isConfirmed} & Представя идентичността на потребителя и служи като входна точка към фирми, количка и история \\
\texttt{Company} & \texttt{name}, \texttt{mol}, \texttt{vat}, \texttt{address}, \texttt{isRegistered} & Представя фирмен субект, който се асоциира с потребителя и към който принадлежат продукти \\
\texttt{Product} & \texttt{name}, \texttt{price}, \texttt{quantity}, \texttt{url}, \texttt{category} & Представя каталожен артикул, който може да бъде визуализиран, филтриран и закупуван \\
\texttt{ShoppingCartEntity} & \texttt{wishedQuantityForOrder}, връзка към продукт и потребител & Представя временно потребителско намерение за покупка \\
\texttt{HistoryEntity} & \texttt{address}, \texttt{date}, \texttt{price}, връзка към потребител & Представя обобщение на финализирана поръчка \\
\bottomrule
\end{longtable}

Таблица \ref{tab:entity-attributes} потвърждава, че домейн моделът е достатъчно компактен, но и достатъчно изразителен. В него няма излишни слоеве на абстракция, а всеки entity клас има ясно разпознаваема роля.

\subsubsection{Изброим тип за категориите}

Категоризационният модел е реализиран чрез \texttt{CategoryEnum}. В него са дефинирани стойности като \texttt{VIDEOCARDS}, \texttt{PROCESSORS}, \texttt{HEADPHONES}, \texttt{PRINTERS}, \texttt{OFFICE\_CHAIRS}, \texttt{SCANNERS}, \texttt{ROUTERS}, \texttt{MICE} и \texttt{KEYBOARDS}. Този избор е особено подходящ по три причини.

Първо, елиминира се рискът от невалидни категории, въведени чрез свободен текст. Второ, контролерът за каталог може да използва ясни branch условия по категории. Трето, категориите се превръщат в част от самия домейн модел, а не остават външен конфигурационен списък без типова защита.

Цената на това решение е, че добавянето на нова категория изисква промяна в кода и повторна компилация. За текущия проект това не е проблем, тъй като категориите са предварително известни и тяхната стабилност е по-важна от динамичната конфигурируемост.

\subsubsection{Binding модели и входна валидация}

Съществена част от домейн обработката в системата не се извършва директно чрез entity класовете, а чрез binding модели. Това са класове, предназначени за приемане на данни от HTML форми. Сред тях са \texttt{UserRegisterBindingModel}, \texttt{UserLoginBindingModel}, \texttt{ProfileUpdateBindingModel}, \texttt{AddCompanyBindingModel}, \texttt{ProductAddBindingModel}, \texttt{AddToCartBindingModel} и \texttt{CheckoutBindingModel}.

От архитектурна гледна точка тези класове изпълняват две роли. Първо, дефинират структурата на потребителския вход. Второ, концентрират на едно място част от валидационната логика чрез анотации като \texttt{@Size}, \texttt{@NotBlank}, \texttt{@NotNull} и \texttt{@Positive} \cite{hibernatevalidator}. Това прави формите по-прогнозируеми и намалява вероятността до service слоя да достигнат очевидно невалидни данни.

\begin{table}[H]
    \centering
    \caption{Примери за валидационни ограничения в binding моделите}
    \label{tab:binding-validation}
    \begin{tabularx}{\textwidth}{L{4.1cm}L{4cm}Y}
        \toprule
        Модел & Ограничения & Смисъл на ограничението \\
        \midrule
        \texttt{ProfileUpdateBindingModel} & \texttt{@Size(min=3, max=20)} за име, фамилия, имейл и пароли & Ограничаване на очевидно невалидни или прекомерно кратки стойности \\
        \texttt{AddCompanyBindingModel} & \texttt{@NotBlank}, \texttt{@Size}, \texttt{@Positive}, \texttt{@NotNull} & Валидиране на фирмено име, адрес и числови идентификатори \\
        \texttt{ProductAddBindingModel} & \texttt{@NotBlank}, \texttt{@Positive}, \texttt{@NotNull} & Гарантира коректни данни за продукт преди запис \\
        \texttt{CheckoutBindingModel} & \texttt{@Size(min=5, max=20)} за адрес & Базова проверка на адреса при финализиране на поръчка \\
        \bottomrule
    \end{tabularx}
\end{table}

Макар този модел да е полезен, следва да се отбележи, че не всички binding класове в проекта имат еднакво богата валидация. Например регистрационният и login моделът са по-леки и не носят симетричен набор от ограничения. Това е напълно допустимо в учебна разработка, но би могло да бъде разширено при бъдещо развитие.

\subsubsection{Service и view модели}

Освен entity и binding класове проектът използва и service модели, и view модели. Service моделите служат като вътрешно представяне между persistence и service слоя, а view моделите са оптимизирани за визуализация в шаблоните или за връщане през API. Примери за такива класове са \texttt{UserServiceModel}, \texttt{UserLoginServiceModel}, \texttt{CompanyServiceModel}, \texttt{CompanyViewModel}, \texttt{ProductViewModel}, \texttt{UserViewModel} и \texttt{MaxPriceViewModel}.

Предимството на този подход е, че отделните слоеве могат да работят с различни представяния на едни и същи логически данни, без да бъдат принудени да използват директно entity обектите навсякъде. Това подобрява капсулацията. Например шаблонът за каталог не е длъжен да получава пълния entity обект на продукта, а само тази част от информацията, която е необходима за визуализация.

\subsubsection{Преобразувания между моделите}

Преобразуването между различните представяния е подпомогнато от \texttt{ModelMapper}. Тази библиотека намалява ръчния код при преобразуване на entity към service модел или на service модел към view модел. В контекста на проекта това улеснява реализацията на списъци с потребители, фирми и продукти. Въпреки това част от домейн логиката остава ръчно реализирана, например когато е необходимо да се създаде изцяло нов entity обект с попълване на конкретни полета и зависимости.

Използването на mapper инструмент е особено полезно при средни по размер проекти. То намалява повторяемия код, но същевременно не скрива домейн решенията, когато те са специфични. Именно такъв е случаят тук: стандартните трансформации се автоматизират, а смислово важните операции се извършват експлицитно в service слоя.

\subsubsection{Съответствие между домейн модел и потребителски сценарии}

Една от силните страни на разглеждания проект е, че домейн моделът сравнително пряко следва основните потребителски сценарии. Потребителят създава фирма; фирмата притежава продукти; продуктът се добавя в количката; количката участва във финализиране на поръчка; поръчката се архивира в историята. Тази линия е достатъчно ясна, за да бъде следвана както от кода, така и от документацията.

Разбира се, моделът не е напълно независим от ограниченията на имплементацията. Историята не съдържа отделни позиции по поръчка, а количката е опростена спрямо типична production система. Въпреки това домейнът е достатъчно консистентен, за да изпълни предназначението си като основа на учебно-приложен онлайн магазин.

\subsubsection{Критични наблюдения върху модела}

Анализът на модела позволява да бъдат формулирани и няколко критични наблюдения:

\begin{itemize}[leftmargin=*]
    \item налице са дребни несъответствия в именуването на някои свойства, което затруднява концептуалната чистота;
    \item поръчката е представена в силно опростен вид чрез исторически запис, без отделни order item обекти;
    \item някои валидационни правила са концентрирани само върху част от binding моделите;
    \item моделът е силно ориентиран към един потребителски тип и не включва роли или права.
\end{itemize}

Тези наблюдения не обезсилват решението, а по-скоро очертават границите на неговата текуща зрялост. Именно поради това те са ценни за документацията: показват, че системата е анализирана критично, а не само описателно.

\subsubsection{Обобщение на проектирането}

Проектирането на уеб-базираната платформа съчетава ясно дефинирани изисквания, добре подредена многослойна архитектура и компактен, но изразителен домейн модел. Тези три аспекта формират основата, върху която стъпва програмната реализация, разгледана в следващата глава.
